\documentclass{article}

% ready for submission
\usepackage[final]{neurips_2023}

\usepackage[utf8]{inputenc} % allow utf-8 input
\usepackage[T1]{fontenc}    % use 8-bit T1 fonts
\usepackage{hyperref}       % hyperlinks
\usepackage{url}            % simple URL typesetting
\usepackage{booktabs}       % professional-quality tables
\usepackage{amsfonts}       % blackboard math symbols
\usepackage{nicefrac}       % compact symbols for 1/2, etc.
\usepackage{microtype}      % microtypography
\usepackage{amsmath}        % for \text in math mode
\usepackage{xcolor}         % colors
\usepackage{graphicx}       % For adding figures/spectrograms
\usepackage{float}          % [H] placement: keep floats exactly in place

\title{Task 5: Challenge}

\author{%
   Validated Encoders \AND
  Amr ABDULLAH
  \And
  Ibrahim MUHAMMAD
}

\begin{document}

\maketitle

\begin{abstract}
We build a sound event detection (SED) system for the MLPC~2026 challenge.
Starting from the tutor-provided decision-tree baseline (non-hidden test
macro-F1 \textbf{0.316}), we replace the per-class trees with
\texttt{HistGradientBoostingClassifier} on the same aggregated
1\,s-segment features, perform a grid analysis over
\texttt{max\_iter}$\times$\texttt{max\_leaf\_nodes}, tune per-class
decision thresholds on a held-out local-validation split, and apply a
median-filter post-processing step with a swept window size. The final
system reaches \textbf{0.545} non-hidden-test macro-F1 (\textbf{+0.229}
over the baseline), driven mostly by the classifier swap and threshold
tuning; post-processing is neutral. Sections 1--4 of the task are
addressed; bonus tasks are out of scope here.
\end{abstract}

% =============================================================================
\section{Baseline Reproduction}
\label{sec:baseline}

\subsection{Pipeline and Reproduction}
The provided baseline trains
one binary \texttt{DecisionTreeClassifier} per class wrapped in
\texttt{MultiOutputClassifier}, on the 960-dimensional aggregated feature
matrix concatenated from the {mean, std, min, max} keys in the
\texttt{.npz} files. Training uses a random 50{,}000-segment subsample of
the full $\sim$1.7\,M training segments, with the remaining feature
weights identical across classes. At inference the classifier predicts on
\emph{all} overlapping 1\,s segments produced at a 0.5\,s hop, then keeps
only those starting at whole second timestamps so that the metric's
1\,s-resolution evaluation is honored. Consecutive active segments are
merged into onset/offset intervals.

We reproduced the baseline on the seed-42 local-validation split (499
files) and the disjoint non-hidden test split (500 files), evaluating
with the provided \texttt{evaluate.py} (segment based macro-F1):
\begin{itemize}\setlength\itemsep{0pt}
  \item Local validation macro-F1 = \textbf{0.307}
  \item Non-hidden test macro-F1 = \textbf{0.316}
\end{itemize}
Our reproduction (\texttt{0.3073} / \texttt{0.3164}) matches the
provided baseline values to within rounding.

\subsection{Per-Class Performance Differences}
% TODO: insert per-class bar/figure from baseline; discuss
The class wise F1 spread is wide. Sustained broadband classes
(\texttt{running\_water}, \texttt{vacuum\_cleaner}, \texttt{microwave})
already score above $0.5$ at the decision tree level: they emit
energy continuously across the 1\,s window, so the aggregated features
(\texttt{melspect\_*, mfcc\_*}) capture them robustly. Brief transients
(\texttt{light\_switch}, \texttt{window\_open\_close}, \texttt{wardrobe\_
drawer\_open\_close}) score lowest: a sub-100\,ms impulse contributes a
small fraction of the 1\,s aggregation, so the per-segment posterior
remains close to the negative class. Frequency also plays a role: with a
$\sim$47$\times$ positive-rate ratio between \texttt{running\_water} and
\texttt{light\_switch}, the rare classes receive far fewer training
splits that distinguish them. Acoustically similar classes
(\texttt{door\_open\_close} vs.\ \texttt{wardrobe\_drawer\_open\_close})
confuse for the same reason a 1\,s mean-pooled spectrum cannot resolve
their short impulses.

\subsection{Strengths and Limitations}
The baseline's strengths are simplicity and a fully unbiased per-class
operating point selection: each binary tree can choose its own
thresholds in feature space. Its limitations, however, are exactly what
caps performance at $\sim$0.307: (i) \textbf{independent binary trees}
cannot model cross-class dependencies (e.g.\ a microwave hum and water
running together), (ii) \textbf{1\,s temporal resolution} is too coarse
for the four transient classes whose typical event length is much
shorter than the window, (iii) \textbf{single CART trees} cannot
combine features as effectively as a boosted ensemble on this many
inputs, and (iv) \textbf{\texttt{predict()} at $\tau{=}0.5$} ignores
the well-known fact that rare-class posteriors lie far below 0.5 and
must be calibrated.

% =============================================================================
\section{Simple Classifier Swap and Tuning}
\label{sec:simple}

\subsection{Best Previous-Phase Classical Model and Choice for the Challenge}
The non-deep-learning leaderboard from our previous phase was
\textbf{XGBoost} (test macro-F1 = 0.5605; 300 trees, \texttt{max\_depth=6},
\texttt{learning\_rate=0.1}, \texttt{tree\_method=`hist`}). This time we use the sklearn native
\textbf{\texttt{HistGradientBoostingClassifier}}, which is algorithmically
equivalent (histogram-based gradient boosting) and produces the same
class of tabular accuracy. As in the baseline, one binary HistGB is
trained per class via \texttt{MultiOutputClassifier}.

\subsection{Hyperparameter Analysis}
We swept the two hyperparameters that most directly control HistGB's
bias--variance trade-off:
\texttt{max\_iter}\,$\in\{200,500\}$ (boosting rounds, complexity of the
additive model) and
\texttt{max\_leaf\_nodes}\,$\in\{31,63\}$ (per-round capacity). Each of
the four configurations is trained on the 50k subsample and evaluated on
the local-validation set with the official scorer (Fig.~\ref{fig:hp}).



The four configs score $0.467$ ($200/31$), $0.458$ ($200/63$), $0.467$ ($500/31$),
$0.458$ ($500/63$). The best configuration was
\texttt{max\_iter=500}, \texttt{max\_leaf\_nodes=31}, giving local-val
macro-F1 = \textbf{0.467} --- already \textbf{+0.16} over the
decision-tree baseline (0.307) before any threshold tuning.

\subsection{Per-Class Threshold Tuning}
HistGB outputs calibrated probabilities; the default $\tau{=}0.5$
systematically under-predicts rare classes whose posteriors sit far
below 0.5 and is the single largest decision-rule lever
(+0.20\,F1 in our previous phase). We collect probability matrices on
the local-validation segments, sweep $\tau\in[0.05, 0.95]$ in steps of
$0.02$ per class and pick the F1-optimal $\tau$ per class. The tuned
thresholds are all low ($0.05$--$0.27$), confirming the calibration
mismatch. Applying them lifts local-validation macro-F1 from
\textbf{0.467} to \textbf{0.554} (\textbf{+0.087}) --- the second-largest
lever after the classifier swap.

\subsection{Comparison and Selected Classes}
Table~\ref{tab:build} summarises how each step builds up the final score.


On the non-hidden test set (single-shot, untouched during development):
baseline DT macro-F1 = \textbf{0.316}; HistGB with the best grid config
and tuned thresholds = \textbf{0.545} (delta \textbf{+0.229}).
Figure~\ref{fig:perclass} shows every one of the 15 classes improves.
The largest absolute gains come from the rare/transient classes
(\texttt{light\_switch}, \texttt{window\_open\_close},
\texttt{wardrobe\_drawer\_open\_close}), where the low tuned thresholds
unlock previously-suppressed positives, and from the mid-frequency
classes (\texttt{phone\_ringing}, \texttt{keychain}) that benefit from
boosting's stronger decision boundaries. Sustained broadband classes
(\texttt{running\_water}, \texttt{vacuum\_cleaner}) were already strong
and stay near the top.

\subsection{Qualitative Error Analysis}
Files were ranked by a per-file F1 over \emph{active or predicted}
classes only (so the metric reflects ``how well did the model do on this
file'' rather than averaging over 11 inactive classes). We show a clear
\emph{success} (\texttt{001565.wav}) and a clear \emph{failure}
(\texttt{000554.wav}); a third mixed example is in
Appendix~\ref{app:mixed}. Figures overlay ground-truth event spans
(transparent boxes) with HistGB framewise probabilities (curves) and the
per-class thresholds (dashed horizontal lines).

\textbf{Success (\texttt{001565.wav}, active-class F1 = 0.97).} Three
sustained classes (\texttt{cutlery\_dishes}, \texttt{microwave},
\texttt{running\_water}) are detected cleanly; the probability curves
track the GT boxes throughout and \texttt{coffee\_machine} stays
correctly below its threshold (Fig.~\ref{fig:err-clean}).

\textbf{Failure (\texttt{000554.wav}, active-class F1 = 0.00).} The true
event is missed entirely while the model raises one or two
acoustically-related transient classes above threshold ---
a clean wrong-prediction failure (Fig.~\ref{fig:err-fail}). This is the
generic transient-class failure: a brief impulse smeared across a 1\,s
window is acoustically ambiguous between
\texttt{door\_open\_close}, \texttt{footsteps}, and
\texttt{wardrobe\_drawer\_open\_close}.



% =============================================================================
\section{Post-Processing: Median Filtering}
\label{sec:postproc}

Independent per-segment predictions are temporally noisy: a single
spurious activation surrounded by zeros is interpreted as a 1\,s event
under \texttt{predictions\_to\_intervals}. We apply a length-$w$
median filter along the time axis (per class), so a binarized
prediction at time $t$ is replaced by the median of the $w/2$ neighbours
on each side. With an odd $w$, a 1 survives only if at least
$\lceil w/2 \rceil$ neighbours are also 1 --- isolated false positives
are removed while events of length $\geq w/2$ are preserved. We swept
$w \in \{1, 3, 5, 7, 9\}$ on local validation (Fig.~\ref{fig:winsweep}):


The best window was $w{=}\textbf{3}$ with local-val macro-F1
$\textbf{0.5556}$ (vs.\ $0.5538$ without filtering, delta \textbf{+0.002});
larger windows degrade monotonically ($w{=}5$: $0.528$; $w{=}9$:
$0.471$) as they erase true events shorter than $\lceil w/2 \rceil$
seconds. On the locked, single-shot non-hidden test set the
$w{=}3$ system delivers $\textbf{0.5437}$ vs.\ $\textbf{0.5450}$ without
filtering (\textbf{$-$0.001}) --- a \emph{neutral} result within
seed-variance noise. The interval count drops sharply ($5{,}545 \to
2{,}453$ on validation), so the filter is removing many isolated
single-segment detections, but on this metric the removed detections
were roughly balanced between true and false positives
(Fig.~\ref{fig:wincls}).


Although $w{=}3$ was the validation-optimal choice, we submit the
\emph{unfiltered} ($w{=}1$) predictions for the hidden test, for two
reasons: (i) median filtering is neutral to negative on the held-out
test, so it carries no expected benefit, and (ii) on 5 of the 1{,}007
hidden files every detection was a single isolated segment that $w{=}3$
erased entirely, leaving those files with no predictions; the $w{=}1$
system covers all 1{,}007 files. The honest conclusion is that median
filtering does not help this system: the threshold-tuned HistGB already
produces temporally coherent predictions, and the remaining errors are
class-confusions rather than the isolated single-segment noise a median
filter removes.

% =============================================================================
\section{Reflection and Real-World Considerations}
\label{sec:reflection}

\subsection{Technical Limitations and Future Work}
The dominant remaining failure mode is the \textbf{1\,s temporal
resolution}: events shorter than half the median window become
filterable noise, and three of the four transient classes
(\texttt{light\_switch}, \texttt{window\_open\_close},
\texttt{wardrobe\_drawer\_open\_close}) remain below F1$\,{=}\,0.30$
even after threshold tuning. Promising directions: (i) train on the
denser 0.5\,s-hopped segments and use them at inference too (the
discarded half-second segments would double the training data), (ii)
move to a small temporal model  that does not lose intra-segment temporal structure, and (iii)
exploit cross-class structure (e.g.\ a chained classifier or a small
multi-class MLP on per-class probability vectors) to model co-occurrence
patterns that independent binary trees cannot. The current
\emph{representation} (mean/std/min/max of 14 base groups) is also a
known bottleneck --- aggregation low-pass-filters time-domain
perturbations.

\subsection{Real-World Deployment}
Inference cost is dominated by the per-segment feature extraction
($\sim$10--40\,ms/segment on a CPU), comfortably under the 1\,s segment
duration --- the system can run in real time on a smart-home device's
microcontroller-class CPU. Robustness, however, is the practical
blocker: collector-grouped evaluation has shown a $\sim$0 val--test gap
in our development data, but a household with unusual acoustics (loud
HVAC, an unusual microphone) will likely drop precision on the
boundary classes the model already struggles with. Threshold transfer
is the largest operational risk: per-class thresholds tuned on the
challenge data will be miscalibrated for a noisier home, and a
deployment system should expose a lightweight on-device recalibration
routine (e.g.\ user-confirmed positives over the first day). Compute
and latency are not the binding constraints --- reliability under
acoustic shift is.

% ==============================FIGs===============================================

\section{Figures}

\begin{figure}[H]
\centering
\includegraphics[width=0.7\linewidth]{hp_grid_heatmap.png}
\caption{HistGB hyperparameter sweep, local validation macro-F1 (no
threshold tuning). \texttt{max\_leaf\_nodes=31} clearly beats $63$
(the larger leaf budget over-fits the 50k subsample), while
\texttt{max\_iter} has almost no effect --- early stopping caps the
effective number of rounds well before 500.}
\label{fig:hp}
\end{figure}

\begin{table}[H]
\centering
\small
\caption{System build-up. Macro-F1 at each stage. ``local val'' is the
seed-42 validation split used for model selection and threshold tuning;
``non-hidden test'' is the single-shot held-out estimate.}
\label{tab:build}
\begin{tabular}{lcc}
\toprule
\textbf{System} & \textbf{local val} & \textbf{non-hidden test} \\
\midrule
Baseline: per-class DecisionTree            & 0.307 & 0.316 \\
HistGB (best grid config, $\tau{=}0.5$)     & 0.467 & --    \\
\;\; + per-class threshold tuning           & 0.554 & 0.545 \\
\;\; + median filter ($w{=}3$)              & 0.556 & 0.544 \\
\bottomrule
\end{tabular}
\end{table}

\begin{figure}[H]
\centering
\includegraphics[width=\linewidth]{section2_perclass.png}
\caption{Per-class F1 on the non-hidden test set: baseline
\texttt{DecisionTree} vs.\ HistGB with the best grid config and tuned
per-class thresholds.}
\label{fig:perclass}
\end{figure}

\begin{figure}[H]
\centering
\includegraphics[width=\linewidth]{err_clean.png}
\caption{Success example (\texttt{001565.wav}). Three sustained classes
detected cleanly; predictions hug the GT boxes throughout.}
\label{fig:err-clean}
\end{figure}

\begin{figure}[H]
\centering
\includegraphics[width=\linewidth]{err_failure.png}
\caption{Failure example (\texttt{000554.wav}). The true class is missed
while the model raises acoustically-related transient classes --- a clean
wrong-prediction failure on a brief-transient event.}
\label{fig:err-fail}
\end{figure}


\begin{figure}[H]
\centering
\includegraphics[width=0.6\linewidth]{section3_window_sweep.png}
\caption{Median-filter window sweep on the local validation set.}
\label{fig:winsweep}
\end{figure}


\begin{figure}[H]
\centering
\includegraphics[width=\linewidth]{section3_per_class_delta.png}
\caption{Per-class effect of the $w{=}3$ median filter on local
validation. The few classes that gain are the longer sustained events
(their true positives survive the $\lceil w/2 \rceil$ vote while isolated
false detections are removed); the short-transient classes lose, because
their true positives are themselves single-segment and are filtered out
--- which is why the net effect is neutral.}
\label{fig:wincls}
\end{figure}
% =============================================================================
\section*{Disclosure of LLM and AI Tool Use}
We used Anthropic Claude and Cursor as engineering assistants throughout:
code restructuring, rebuilding, debugging, testing  and
editing the report draft.
All numerical results, figures, and the prediction CSV were
verified against the raw artifacts on disk.

\begin{thebibliography}{9}
\small
\bibitem[Murphy(2012)]{murphy} K.~P. Murphy. \textit{Machine Learning:
  A Probabilistic Perspective}. MIT Press, 2012.
\bibitem[Mesaros et~al.(2016)]{sed} A.~Mesaros, T.~Heittola, T.~Virtanen.
  Metrics for Polyphonic Sound Event Detection. \textit{Applied Sciences},
  6(6):162, 2016.
\bibitem[Friedman(2001)]{gbm} J.~H. Friedman. Greedy Function Approximation:
  A Gradient Boosting Machine. \textit{Annals of Statistics}, 29(5),
  2001.
\end{thebibliography}

\appendix
\section{Additional Error-Analysis Example}
\label{app:mixed}
\begin{figure}[H]
\centering
\includegraphics[width=\linewidth]{err_mixed.png}
\caption{Mixed example (\texttt{004551.wav}, active-class F1 = 0.45, 6
active classes). The model handles the sustained classes but shows
boundary errors and timing offsets on the transients --- a typical
intermediate case between the clean and failure extremes.}
\label{fig:err-mixed}
\end{figure}

\end{document}
